\documentclass[UTF8,a4paper,11pt]{ctexart}

\usepackage[margin=2.35cm]{geometry}
\usepackage{amsmath,amssymb,bm}
\usepackage{booktabs}
\usepackage{threeparttable}
\usepackage{tabularx}
\usepackage{array}
\usepackage{graphicx}
\usepackage{float}
\usepackage{enumitem}
\usepackage[numbers,sort&compress]{natbib}
\usepackage{xcolor}
\usepackage{microtype}
\usepackage{hyperref}

\hypersetup{
  colorlinks=true,
  linkcolor=blue!55!black,
  citecolor=green!35!black,
  urlcolor=blue!65!black
}
\setlength{\parindent}{2em}
\setlength{\parskip}{0.25em}
\setlist{nosep,leftmargin=2em}
\newcolumntype{Y}{>{\centering\arraybackslash}X}

\newcommand{\zmethod}{Z0}
\newcommand{\gmethod}{G4}
\newcommand{\placeholderfigure}[2]{%
  \begin{figure}[H]
    \centering
    \fbox{\parbox[c][4.2cm][c]{0.88\linewidth}{\centering\textbf{图待补}\par #1}}
    \caption{#2}
  \end{figure}
}

\title{历史 OD 驱动的选择性区域物化：\\
面向精确最短路查询的无参数需求传播与神经排序残差}
\author{席圣洋（第一作者） \and 肖逸涵（第二作者） \and 王加文（第二作者）}
\date{初稿日期：2026 年 8 月 1 日}

\begin{document}

\maketitle

\begin{center}
作者单位：\rule{0.66\linewidth}{0.4pt}\\[0.6em]
通讯作者：\rule{0.22\linewidth}{0.4pt}\qquad
电子邮箱：\rule{0.30\linewidth}{0.4pt}
\end{center}

\begin{abstract}
道路网络分析和交通仿真需要反复执行大规模起终点（origin--destination, OD）最短路查询。
全图在线搜索开销较高，而全量索引预处理又会占用大量存储。本文在已有 CRP 式边界覆盖图精确
索引之上，研究一个更具体的问题：如何仅根据历史 OD 与静态道路图，在固定预算下选择少量值得
物化的连通区域，同时保持在线查询距离完全精确。为此，本文构造全图空间均匀候选和基于真实
查询展开节点减少量的精确区域收益标签，并采用历史、标签和未来三个互不重叠的连续时间窗口，
以及基于 Jaccard 重叠连通组的空间隔离。本文首先提出面向该任务的无参数方法 \zmethod：分别将
历史起点和终点边际需求沿有向道路图正、反方向进行固定 lazy 扩散，在多尺度上聚合区域需求暴露。
在此基础上，\gmethod 采用 32 层双塔图传播网络学习冻结 \zmethod 无法解释的排序残差，并以
可微 Spearman 目标和仅依赖验证集的峰值保存协议稳定训练。Porto 与 Chicago 两城实验表明，
\gmethod 的三个种子在空间留出集和未来窗口的全局 Spearman 上均超过各自 \zmethod；但这种
全局排序增益并不保证所有 Top-$K$ 或在线收益同步提高。严格非重叠部署下，两城全部在线查询均
保持 100\% 距离正确。C++ 同实现评测进一步表明，Porto 的平均查询耗时下降 6.51\%--7.48\%，
而 Chicago 的收益主要体现为展开节点和多数 P95 尾延迟下降，平均耗时受密集 shortcut 扫边开销
限制。结果说明，历史需求传播能够为精确最短路索引提供稳健、可解释的区域资源配置依据，而
神经残差更适合作为全局排序增强，而非无条件替代确定性方法。
\end{abstract}

\noindent\textbf{关键词：}道路网络；最短路索引；历史 OD；选择性区域物化；图扩散；学习排序

\section{引言}

最短路距离查询是交通仿真、出行需求分析、位置服务和道路网络研究中的基础操作。经典
Dijkstra 算法能够给出精确答案，但当同一城市路网上持续到来大量 OD 查询时，重复搜索会造成
显著计算开销。已有研究通过层次化预处理、标签索引或区域覆盖图减少在线工作量
\citep{dijkstra1959,salt2014}。这些方法通常面向整张图构建索引，而真实查询负载往往具有明显
的空间偏斜和时间局部性：有限预算下，并非所有道路区域都值得获得相同的预处理资源。

负载感知最短路研究已经证明，历史查询分布可用于调整精确索引
\citep{wan2022wcf,dan2025fahl}；查询缓存工作也利用历史日志估计可复用对象的价值
\citep{thomsen2012cache,li2020cache}。然而，标签索引、路径缓存和区域边界覆盖图对应不同的
物化对象与成本结构。本文不重新发明底层边界 shortcut，而是把问题收窄到：在 CRP 式精确覆盖图
底座上，只物化少量候选区域时，应如何根据历史 OD 对候选进行排序，并在时空隔离、区域互斥和
真实在线查询约束下验证该排序。

这一问题具有三个困难。第一，若用历史最短路径直接统计区域命中，会把昂贵的路径计算提前放入
输入阶段，也容易使方法退化为路径缓存。第二，候选区域之间可能高度重叠；随机划分会让近重复
区域跨越训练集和测试集，造成严重空间泄漏。第三，离线相关性、单区域收益和多区域在线加速并不
等价：重叠候选无法同时物化，shortcut 数量和端点接入还会改变实际查询成本。

本文将 AI 限定为离线资源配置工具。模型不预测最短路，也不改变距离答案；在线阶段仍由精确
双向 Dijkstra 和区域局部接入完成查询。本文的主要贡献如下：

\begin{enumerate}
  \item 在既有负载感知精确索引和 CRP 式边界覆盖图基础上，系统定义历史 OD 驱动的选择性区域
  物化任务，以真实展开节点减少量作为候选区域价值，并保持学习选区与精确查询解耦。
  \item 构建无历史路径输入的无参数双向多尺度需求传播组合 \zmethod，并通过两城的方向、拓扑、
  需求来源、深度和池化正交消融解释其有效信号来源。
  \item 构建以 \zmethod 为冻结先验的 32 层神经排序残差 \gmethod。该方法采用稳定的严格恒等
  残差路径和与评价一致的 soft Spearman 目标，在 Porto 与 Chicago 的多种子空间留出集和未来
  窗口上稳定改善全局排序。
  \item 建立 H/Y/F 时间隔离、Jaccard 空间重叠组隔离、严格非重叠部署和精确在线配对的验证闭环，
  并通过 C++ 同实现评测揭示不同城市中展开节点、shortcut 密度、平均延迟和尾延迟之间的差异。
\end{enumerate}

围绕上述贡献，实验集中回答五个研究问题：RQ1，历史需求传播能否在空间隔离候选上预测区域
价值；RQ2，冻结排序能否推广到未参与建模的未来窗口；RQ3，\zmethod 的收益究竟来自历史需求、
正确道路拓扑、传播方向还是简单池化统计；RQ4，单区域排序能否在严格不重叠的多区域部署中转化
为精确查询工作量收益；RQ5，同实现下的工作量减少能否进一步转化为真实墙钟加速。本文按此顺序
组织实验，避免用单一相关系数替代系统结论。

\placeholderfigure{方法总图：H/Y/F 时间轴、候选区域、\zmethod、\gmethod、hard-disjoint 选择、
CRP 式覆盖图和精确查询。}{本文端到端流程。底层覆盖图采用已有精确索引结构，本文研究重点是
历史 OD 驱动的区域排序、隔离与选择性物化。}

\section{相关工作}

\subsection{精确最短路索引与区域覆盖图}

道路最短路加速包含层次化搜索、标签索引和基于分区的覆盖图等路线。SALT 将图分区、边界覆盖图
与查询算法组合到统一框架中\citep{salt2014}。其中，对区域内部运行受限最短路、在边界节点之间
建立精确 shortcut、删除内部节点后在覆盖图上查询，与本文的底层压缩结构高度一致。因此本文
明确把该部分视为已有索引底座，不将其列为原创算法。本文的研究对象是只选择少量区域物化时的
候选排序、内部端点接入、严格不重叠部署和同预算在线效果。

\subsection{负载感知索引与最短路缓存}

WCF 利用道路查询负载的空间偏斜与时间局部性构造 workload-aware Core-Forest 标签索引
\citep{wan2022wcf}；FAHL 进一步研究 flow-aware 的标签索引\citep{dan2025fahl}。另一些工作从
历史查询日志中选择值得缓存的路径或区域对象\citep{thomsen2012cache,li2020cache}。这些工作说明
“历史负载驱动精确最短路加速”并非新问题。与之不同，本文不预测高频节点后构建 2-hop/Core-
Forest 标签，也不缓存已出现的完整路径，而是对固定、空间均匀的连通区域预测其精确展开收益，
并在 CRP 式覆盖图中进行选择性物化。

\subsection{图扩散、深层传播与可微排序}

DCRNN 在有向道路图上使用双向随机游走建模空间依赖\citep{li2018dcrnn}；SIGN 和 Jumping
Knowledge 展示了多尺度图扩散和跨层表示读出的价值\citep{frasca2020sign,xu2018jk}。因此，本文
不把 \zmethod 的正反向多步传播写成新的基础扩散算子，而强调其在“起点/终点边际注入--固定
lazy 扩散--区域 mean+max--精确收益排序”上的任务化组合。

\gmethod 的传播组织受到 NBFNet 的 Indicator--Message--Aggregate 框架启发
\citep{zhu2021nbfnet}。深层网络中的 pre-activation 残差和小尺度残差已有明确先例
\citep{li2020deepergcn,touvron2021layerscale}，可微排序亦是成熟方向\citep{jagerman2022rax}。
本文不将这些组件分别作为新通用架构，而研究它们能否在强确定性先验上形成稳定、可审计的任务
增量。

\section{问题定义}

\subsection{道路图、查询负载与时间窗口}

给定带正权有向道路图 $G=(V,E,w)$，其中 $w:E\rightarrow\mathbb{R}_{>0}$ 表示道路长度。
一个查询 $q=(o_q,d_q,t_q)$ 包含起点、终点和时间戳。全部查询按时间排序后划分为三个连续且
互不重叠的窗口：历史窗口 $H=[0,0.35)$、当前标签窗口 $Y=[0.35,0.70)$ 和未来冻结窗口
$F=[0.70,1.00)$。$H$ 只用于构造需求输入，$Y$ 用于生成训练标签和 train/validation 选型，
$F$ 只在协议冻结后进行零样本时间评测。

历史窗口仅保留起终点边际需求。对节点 $v$，定义
\begin{equation}
  h_O(v)=\frac{1}{|H|}\sum_{q\in H}\mathbf{1}[o_q=v],\qquad
  h_D(v)=\frac{1}{|H|}\sum_{q\in H}\mathbf{1}[d_q=v].
\end{equation}
方法不读取 $H$ 中任一查询的最短路径。

\subsection{候选区域与精确收益标签}

候选集合 $\mathcal{R}=\{R_1,\ldots,R_M\}$ 在全图节点上以固定随机种子尽量空间均匀地产生，
每个候选是包含 512 个节点的连通子图。候选生成不读取历史 OD、标签或模型输出。区域 $R$ 的
边界集合 $B_R$ 包含与区域外节点相邻的节点，其余节点构成内部集合 $I_R$。

对标签查询 $q$，记原图双向 Dijkstra 的展开节点数为 $N_G(q)$，只物化区域 $R$ 后的覆盖图搜索
与端点局部接入总展开节点数为 $N_R(q)$。区域的监督标签定义为
\begin{equation}
  y_R^{Y}=\frac{1}{|Q_Y|}\sum_{q\in Q_Y}\left[N_G(q)-N_R(q)\right],
  \label{eq:label}
\end{equation}
其中 $Q_Y$ 是从 $Y$ 中以预先固定随机种子抽取的 2,000 条查询。未来标签 $y_R^F$ 使用完全
相同定义，但仅在最终评测时从 $F$ 生成。标签衡量精确算法工作量而非不稳定的单次 Python
墙钟时间；每条查询同时校验索引距离与原图距离相等。

\subsection{选择性物化目标}

排序器输出候选分数 $s_R$。部署阶段按分数从高到低贪心选择 $K\in\{5,10,18\}$ 个区域，且要求
任意两个已选区域节点集合严格不相交：
\begin{equation}
  |S|=K,\qquad R_i\cap R_j=\varnothing,\quad \forall R_i\ne R_j\in S.
\end{equation}
这一 hard-disjoint 约束来自当前索引的物化语义。式\eqref{eq:label}是单区域边际收益，多区域
收益不能简单相加，因此最终价值必须通过将整个 $S$ 一次性物化后的真实在线查询确定。

\section{方法}

\subsection{CRP 式选择性区域物化}

对每个被选择区域 $R$，离线阶段从每个 $u\in B_R$ 出发，在诱导子图 $G[R]$ 上运行受限
Dijkstra，并对每个可达的 $v\in B_R$ 建立有向 shortcut：
\begin{equation}
  w_R(u,v)=d_{G[R]}(u,v).
\end{equation}
随后删除 $I_R$，保留外部节点、边界节点、剩余原始边和 shortcut。该结构沿用已有 CRP/边界
覆盖图思想\citep{salt2014}。若查询端点位于某个 $I_R$，算法只在该区域内部计算端点至边界的
精确距离，并将这些距离作为双向搜索的初始 frontier；若起终点位于同一区域，还比较区域内部
直达距离。由此避免整条查询回退原图，同时保持精确性。

\paragraph{距离精确性。}
任意穿越区域 $R$ 的原图路径都可按进入和离开 $R$ 的边界节点切分。将其中每段区域内子路径替换
为对应 shortcut，不会增加路径长度，因为 shortcut 保存的是 $G[R]$ 内的最短距离；反过来，每条
shortcut 都对应 $G[R]$ 中一条真实路径，因而覆盖图不可能产生比原图更短的非法距离。内部端点
通过区域内精确 frontier 接入，起终点同区时另行比较区内直达路径。因此在已选区域互不重叠、
shortcut 精确且边权非负的条件下，覆盖图查询与原图最短路距离相等。实现仍逐查询进行数值校验，
不以理论论证替代工程正确性测试。

\paragraph{预处理与查询成本。}
设区域 $R$ 的边界数为 $b_R$、内部诱导边数为 $m_R$。直接从每个边界节点运行一次受限 Dijkstra，
预处理复杂度约为 $O\!\left(b_R(m_R+|R|)\log |R|\right)$，潜在 shortcut 数为 $O(b_R^2)$。
因此相同的 512 节点预算并不意味着相同的索引成本：边界更稠密的区域会产生更多 shortcut 和在线
扫边。本文同时报告区域节点预算、shortcut 数、展开节点与扫描边，正是为了避免把节点数相同
误认为系统成本相同。

\subsection{无参数双向多尺度需求场 \zmethod}

首先对历史边际需求进行对数缩放：
\begin{equation}
  \bm{o}^{(0)}=\log(1+1000\bm{h}_O),\qquad
  \bm{d}^{(0)}=\log(1+1000\bm{h}_D).
\end{equation}
令 $P_{\rightarrow}$ 表示沿原图方向、按接收节点归一化的均值传播矩阵，$P_{\leftarrow}$ 表示在
反图上的对应传播矩阵。每一层采用带自环的固定 lazy mean：
\begin{equation}
  \bm{o}^{(k+1)}=\tfrac{1}{2}(I+P_{\rightarrow})\bm{o}^{(k)},\qquad
  \bm{d}^{(k+1)}=\tfrac{1}{2}(I+P_{\leftarrow})\bm{d}^{(k)}.
  \label{eq:z0diffusion}
\end{equation}
在深度集合 $\mathcal{D}=\{1,2,4,8,16,32\}$ 上，区域分数为
\begin{equation}
  z_R=\frac{1}{|\mathcal{D}|}\sum_{k\in\mathcal{D}}
  \left[
  \operatorname{mean}_{v\in R} e_v^{(k)}+
  \operatorname{max}_{v\in R} e_v^{(k)}
  \right],\quad
  \bm{e}^{(k)}=\bm{o}^{(k)}+\bm{d}^{(k)}.
  \label{eq:z0score}
\end{equation}
\zmethod 不含可学习参数，不读取 $Y$ 或 $F$。其目标不是预测路径，而是估计候选区域在多尺度
道路邻域中暴露于历史起终点需求的程度。

\subsection{以 \zmethod 为先验的神经排序残差 \gmethod}

\gmethod 包含起点塔和终点塔，二者分别沿正图和反图传播且参数不共享。每层使用边特征门控的
message--aggregate 分支 $\mathcal{F}_\ell$，并采用严格恒等的 pre-norm 小残差更新：
\begin{equation}
  \bm{h}^{(\ell+1)}=\bm{h}^{(\ell)}+
  0.01\,\mathcal{F}_\ell\!\left(\operatorname{LayerNorm}(\bm{h}^{(\ell)}),E\right).
  \label{eq:g3}
\end{equation}
该结构用于抑制旧 32 层主干中的反向指数放大；它是已有深层残差思想在本任务中的稳定化实现，
不是新的通用残差结构。模型在 $1,2,4,8,16,32$ 层读取双塔状态并进行区域 mean+max 池化，输出
残差 $r_\theta(R)$。最终分数为
\begin{equation}
  s_R=z_R+\alpha r_\theta(R),\qquad
  \alpha\in\{0,0.125,0.25,0.5,1\}.
  \label{eq:g4score}
\end{equation}
$\alpha=0$ 严格回退到 \zmethod；每个 epoch 的 $\alpha$ 仅在 validation 上选择，同分时优先更小
残差。

\subsection{与评价一致的 soft Spearman 目标}

旧 pairwise softplus 目标可通过放大分数尺度持续下降，却不保证真实 Spearman 同向变化。为减少
这种目标错位，先将批内分数标准化为 $\hat{s}_i$，再用成对 sigmoid 构造软排名
\citep{jagerman2022rax}：
\begin{equation}
  \widetilde r_i=\sum_j \sigma\!\left(\frac{\hat{s}_i-\hat{s}_j}{\tau}\right),
  \qquad \tau=0.1.
\end{equation}
设真实标签的平均秩为 $r_i^*$，训练损失为
\begin{equation}
  \mathcal{L}_{\mathrm{rank}}=1-\operatorname{Pearson}
  (\widetilde{\bm r},\bm r^*).
\end{equation}
所有模型以 validation Spearman 选择 checkpoint。holdout 和 $F$ 不参与结构、学习率、epoch、
$\alpha$ 或策略选择。

\section{实验设计}

\subsection{数据集与图规模}

\begin{table}[H]
  \centering
  \caption{两城数据与冻结切分。}
  \label{tab:data}
  \resizebox{\textwidth}{!}{%
  \begin{tabular}{lrrrrrrr}
    \toprule
    城市 & 交通数据 & 节点数 & 有向边数 & 全部 OD & $H$ & $Y$ & $F$ \\
    \midrule
    Porto & 出租车 & 133,839 & 221,589 & 98,082 & 34,328 & 34,329 & 29,425 \\
    Chicago & 共享单车 & 108,114 & 200,155 & 100,000 & 35,000 & 35,000 & 30,000 \\
    \bottomrule
  \end{tabular}}
\end{table}

Porto 使用公开出租车轨迹数据\citep{moreira2013porto}，将有效轨迹首末点映射为 OD；Chicago
使用 Divvy 2022 年 6 月共享单车行程\citep{divvydata}。Chicago 数据先按字段合法性、时间顺序、
研究边界和路网吸附质量清洗，再按完整时间序列等间隔抽取 100,000 条，不依据模型或区域收益
筛选。道路网络来自 OpenStreetMap\citep{openstreetmap}；Chicago 固定使用与行程时间接近的
2022 年 1 月 Illinois 历史快照，而不是按实验结果改用最新路网。每城生成 1,200 个 512 节点
候选，并划分为 840/180/180 个 train/validation/holdout 候选。对候选两两计算 Jaccard 重叠，
当重叠不低于 0.50 时连边，再将整个连通组分配到同一 split。Porto 与 Chicago 分别形成 521 和
452 个重叠组，跨 split 不存在 Jaccard 不低于 0.50 的候选对。

每个城市的 $Y$ 与 $F$ 都对全部 1,200 个候选、各 2,000 条固定查询生成标签，因此两城共包含
960 万次候选--查询标签评估。候选、查询抽样种子和 split 在读取 holdout/future 结果前冻结。

\subsection{泄漏控制与冻结顺序}

候选区域是空间对象，相邻 BFS 区域常共享大量节点。早期随机 candidate split 中，直接用训练集
最近重叠区域的标签预测测试候选，Spearman 可达约 0.9936，表明普通随机划分会严重高估泛化。
本文因此先建立阈值 0.50 的 Jaccard 重叠图，再以连通组为最小分配单元。该规则比只删除完全
重复区域更严格，也避免模型通过共享大部分节点“记住”测试区域。

时间上，$H$、$Y$ 和 $F$ 的查询 ID 与时间范围均无交叉。候选身份在读取 $H$ 前由静态道路图
确定；模型输入只来自 $H$；训练与 validation 标签来自 $Y$；holdout 只用于空间外部评测；$F$
在结构、超参数、种子集合和模型选择策略全部冻结后才解锁。Chicago 只复用 Porto 已确定的协议，
不依据其 holdout、future 或在线结果重新调整深度、学习率、温度与残差门。

\subsection{对照方法}

本文比较以下方法：随机排序；按历史端点频率选择的热点方法；只使用区域几何中心位置的 midpoint
Proxy；无区域传播的 MLP；无参数 \zmethod；以及神经残差 \gmethod。MLP 报告 5 个随机种子，
\gmethod 报告 42/43/44 三个种子。所有方法共享相同候选池和部署预算，不允许根据 holdout 或
future 删除不利种子。

其中随机与历史热点回答“是否仅靠任意区域或端点频率即可获益”；midpoint Proxy 检验简单空间
位置是否已经足够；MLP 检验在不传播道路邻接时监督区域特征能够达到的水平；\zmethod 检验固定
图扩散；\gmethod 则只检验监督残差能否超过强先验。Oracle 只用于描述标签空间上界，不作为可
部署方法，也不参与任何超参数选择。

\subsection{指标与实现}

离线排序报告 Spearman、NDCG@$K$ 和 Top-$K$ 平均真实收益，$K\in\{5,10,18\}$。部署阶段报告
shortcut 数、平均展开节点变化、平均耗时、P95 耗时和距离正确率。在线变化率均相对同实现原图
双向 Dijkstra 计算。\gmethod 含 418,183 个参数，采用 CUDA FP32、学习率 $5\times10^{-3}$、
40 个 epoch、无 scheduler。训练只在 GPU 上进行。

C++ 评测在同一 C++20/O3 可执行程序内实现原图与压缩图双向 Dijkstra，固定单 CPU 核，预热
2 轮并正式重复 10 轮。该设计避免把不同语言或不同机器的计时差异误当作算法收益。

神经模型训练在 RTX 4090 D 上使用 CUDA；CPU 只承担数据审计、单元测试和最终 C++ 查询评测。
均值与标准差按预先固定的全部种子报告。本文不把不同种子数的方法比较包装为显著性检验，也不
在观察 holdout 后删除异常种子。在线计时采用逐查询原图--压缩图配对，查询顺序和实现保持一致；
算法无关的系统扰动通过完整干净复跑而不是选择性删除样本处理。

\section{实验结果}

\subsection{RQ1：历史需求传播能否预测区域价值？}

\begin{table}[H]
  \centering
  \caption{空间 holdout 排序结果。MLP 为五种子均值，\gmethod 为三种子均值。}
  \label{tab:holdout}
  \begin{tabular}{llcc}
    \toprule
    城市 & 方法 & Spearman & NDCG@18 \\
    \midrule
    Porto & midpoint Proxy & 0.891756 & \textbf{0.983040} \\
           & MLP & $0.869552\pm0.019732$ & $0.959260\pm0.014273$ \\
           & \zmethod & 0.935603 & 0.974533 \\
           & \gmethod & $\mathbf{0.945902\pm0.001699}$ & $0.974551\pm0.000110$ \\
    \midrule
    Chicago & midpoint Proxy & 0.895597 & 0.867650 \\
            & MLP & $0.847990\pm0.023984$ & $0.901271\pm0.028814$ \\
            & \zmethod & 0.941135 & \textbf{0.922494} \\
            & \gmethod & $\mathbf{0.946787\pm0.001144}$ & $0.866699\pm0.033993$ \\
    \bottomrule
  \end{tabular}
\end{table}

\zmethod 在两城的全局 Spearman 上均高于 midpoint Proxy 与 MLP，说明道路拓扑上的历史需求传播
可以形成强区域价值轴。\gmethod 相对 \zmethod 的 Porto 三种子增量为
$+0.012373/+0.010313/+0.008212$，Chicago 为
$+0.006405/+0.006515/+0.004035$，六次比较方向一致。然而 Chicago 的 NDCG@18 明显下降，
说明 soft Spearman 改善的是全局名次，不是所有头部预算。

\subsection{RQ2：冻结方法能否推广到未来窗口？}

\begin{table}[H]
  \centering
  \caption{冻结未来窗口结果。}
  \label{tab:future}
  \begin{tabular}{llccc}
    \toprule
    城市 & 方法 & Spearman & NDCG@5 & NDCG@18 \\
    \midrule
    Porto & \zmethod & 0.944656 & \textbf{0.944421} & \textbf{0.933698} \\
           & \gmethod & $\mathbf{0.956598\pm0.001580}$ & $0.893761\pm0.005628$ & $0.919343\pm0.001397$ \\
    \midrule
    Chicago & \zmethod & 0.897509 & 0.659842 & \textbf{0.786935} \\
            & \gmethod & $\mathbf{0.925712\pm0.000866}$ & $\mathbf{0.697910\pm0.046648}$ & $0.765981\pm0.030726$ \\
    \bottomrule
  \end{tabular}
\end{table}

Porto 和 Chicago 的三个 \gmethod 种子在 future Spearman 上均超过各自 \zmethod。Porto 的
$Y/F$ 区域收益 Spearman 高达 0.9959，因此该结果证明的是稳定需求时段下的时间外泛化，不能
外推为强分布漂移鲁棒性。Chicago 的时间漂移更明显，\gmethod 的全局增量更大，但 Top-$K$ 仍
存在波动。

\placeholderfigure{两城 holdout/future 的 \zmethod 与 \gmethod Spearman、NDCG@5 和 NDCG@18
分组图，误差线显示随机种子标准差。}{两城排序性能及全局指标与头部指标的差异。}

\subsection{RQ3：\zmethod 的有效信号来自哪里？}

\begin{table}[H]
  \centering
  \caption{正确道路拓扑与度保持随机重连的 Spearman 对照。}
  \label{tab:rewire}
  \begin{tabular}{lccc}
    \toprule
    城市与窗口 & \zmethod & 度保持随机重连 & 绝对下降 \\
    \midrule
    Porto holdout & 0.935603 & 0.556581 & 0.379022 \\
    Porto future & 0.944656 & 0.572467 & 0.372189 \\
    Chicago holdout & 0.941169 & 0.605370 & 0.335799 \\
    Chicago future & 0.897844 & 0.454814 & 0.443030 \\
    \bottomrule
  \end{tabular}
\end{table}

度保持重连保留每个节点的入度和出度，却破坏真实道路邻接；两城四个切分均出现大幅下降，支持
“正确道路拓扑对全局区域价值轴具有独立贡献”。仅起点和仅终点均低于双向 \zmethod，且终点场
贡献更强。无向图的全局相关性接近 \zmethod，但未来 NDCG@5 更低，说明方向信息主要影响头部。
保持两侧边际需求的 OD 重耦合几乎不改变结果，因此现有证据支持边际需求传播，不支持原始 OD
配对关系是必要因素。max pooling 几乎复现全局排序，mean pooling 明显较弱；多深度读出在不同
切分间更稳健，但不存在对所有指标都最优的单一深度。

\subsection{RQ4：排序能否转化为精确在线收益？}

直接按 \zmethod 分数选择 Porto Top-18 时，候选并集只有 2,773 个唯一节点，成员冗余系数为
3.323，并存在共享 116 个节点的候选对，无法由当前索引共同物化。hard-disjoint 贪心最终选满
18 个互不重叠区域，覆盖 9,216 个唯一节点。该变化牺牲部分单区域平均收益，却得到语义明确的
可部署集合，说明“预测高价值区域”和“组成高价值区域集合”是两个不同层次的问题。

\begin{table}[H]
  \centering
  \caption{$K=18$ 严格非重叠部署的在线展开节点结果。负值表示相对原图减少。}
  \label{tab:online}
  \resizebox{\textwidth}{!}{%
  \begin{tabular}{llrrr}
    \toprule
    城市 & 方法 & Shortcuts & 当前展开变化 & 未来展开变化 \\
    \midrule
    Porto & 随机 & 7,986 & $-4.91\%$ & $-4.73\%$ \\
           & 历史热点 & 7,235 & $-15.84\%$ & $-15.90\%$ \\
           & midpoint Proxy & 8,307 & $-17.35\%$ & $-17.66\%$ \\
           & \zmethod & 6,712 & $-17.46\%$ & $-17.76\%$ \\
           & \gmethod（三种子） & 7,755 & $-17.59\%$ & $-17.80\%$ \\
    \midrule
    Chicago & 随机 & 26,011 & $-4.40\%$ & $-4.07\%$ \\
            & 历史热点 & 20,093 & $-17.37\%$ & $-17.95\%$ \\
            & midpoint Proxy & 21,516 & $-20.30\%$ & $-21.08\%$ \\
            & \zmethod & 23,378 & $-21.47\%$ & $-21.92\%$ \\
            & \gmethod（三种子） & 20,718 & $-19.62\%$ & $-20.28\%$ \\
    \bottomrule
  \end{tabular}}
\end{table}

Porto 中，\gmethod 相对 \zmethod 仅多减少约 0.13/0.04 个百分点的当前/未来展开节点，同时
shortcut 增加约 15.5\%，因此 \zmethod 的性能--存储性价比更高。Chicago 中，\zmethod 在展开
节点上优于 \gmethod 与 Proxy。两城既有和新增在线评测合计 204,000 个查询--索引配对，距离
正确率均为 100\%。该结果再次表明，全局排序增量不能直接等同于部署收益。

\subsection{RQ5：算法工作量减少是否转化为墙钟加速？}

\begin{table}[H]
  \centering
  \caption{C++20/O3 同实现的 $K=18$ 性能变化。负值表示降低。}
  \label{tab:cpp}
  \resizebox{\textwidth}{!}{%
  \begin{tabular}{llrrrrrr}
    \toprule
    城市 & 方法 & 当前均值 & 未来均值 & 当前 P95 & 未来 P95 & 当前扫边 & 未来扫边 \\
    \midrule
    Porto & \zmethod & $-6.51\%$ & $-7.48\%$ & $-4.72\%$ & $-6.24\%$ & $-6.67\%$ & $-6.91\%$ \\
           & \gmethod & $-6.41\%$ & $-8.90\%$ & $-5.56\%$ & $-8.96\%$ & $-5.85\%$ & $-6.05\%$ \\
    Chicago & \zmethod & $+3.57\%$ & $+0.72\%$ & $-3.13\%$ & $-3.85\%$ & $+26.56\%$ & $+27.51\%$ \\
            & \gmethod & $+3.03\%$ & $+3.02\%$ & $-2.24\%$ & $-1.33\%$ & $+18.58\%$ & $+19.26\%$ \\
    \bottomrule
  \end{tabular}}
  \vspace{0.35em}

  \parbox{0.94\textwidth}{\footnotesize 注：\gmethod 为三种子均值；表中省略标准差，完整数值保留
  在实验产物中。}
\end{table}

C++ 评测包含 560,000 个正式计时配对样本，距离正确率为 100\%，并精确重放 Python 的区域与
shortcut。Porto 的展开与扫边同时减少，因而平均耗时和 P95 均下降。Chicago 虽减少约 20\% 的
展开节点，但密集边界 shortcut 使扫边增加 18\%--28\%，抵消了典型查询的收益；14 个方法--窗口
组合中仍有 11 个 P95 下降。区域物化是否加速因此取决于边界形状与 shortcut 密度，而不仅取决于
区域排序精度。

\placeholderfigure{Porto 与 Chicago 的展开节点变化、扫边变化、平均耗时和 P95 变化四联图。}
{选择性物化在两城中的系统收益与 shortcut 扫边成本。}

\section{讨论}

\subsection{确定性主方法与神经增强的分工}

实验支持一个分层结论。\zmethod 无需标签和训练，已经获得较强的全局排序与在线工作量收益，
适合作为部署主方法。\gmethod 在两城六个“种子--窗口”组合上稳定提高全局 Spearman，证明神经
网络至少学习到了冻结先验之外的排序信息；但其 Top-$K$ 与在线收益并未全面超过 \zmethod。
因此本文不把神经模型描述为无条件替代传统方法，而将其定位为需要结合预算与头部指标选择的
全局排序增强。

\subsection{从失败机制得到的可复现经验}

旧 32 层网络曾出现约 22 个 epoch 不更新，随后突然变化。逐步诊断发现，FP16 的默认动态缩放
导致早期 optimizer step 被跳过；在 FP32/BF16 下，梯度元素本身仍有限，但全局范数的 FP32
平方和发生溢出，裁剪系数继而接近零。32 层反向 hook 从末层约 $10^{-6}$ 放大到首层约
$10^{19}$。式\eqref{eq:g3}中的严格恒等路径使深层训练稳定。此后又发现旧 pairwise loss 与
Spearman 方向错位，最终才冻结 soft Spearman 协议。该过程说明“数值可训练”和“任务目标有效”
是两个必须分别检验的问题。

\subsection{适用范围与局限性}

本文仍有以下局限。第一，底层边界覆盖图是成熟结构，本文贡献不在发明新的压缩索引。第二，
历史负载驱动精确索引已有 WCF 等直接前作，本文的新意限于选择性区域物化对象、无路径输入、
精确区域收益与完整隔离部署组合。第三，两城实验尚不足以覆盖强需求突变、道路权重实时变化和
不同路网密度；尤其 Porto 的 $Y/F$ 高相关不能代表强漂移。第四，候选固定为 512 节点，标签构造
成本较高，区域大小和边界稀疏化仍需研究。第五，\gmethod 优化全局 Spearman，尚未解决 Top-$K$
和 shortcut 成本的多目标一致性。第六，空间均匀候选只能说明覆盖范围更广；在没有人口、行政区
和可达性公平标签时，不能宣称改善偏远地区公平性。

未来可在不改变本文结论的前提下研究预算感知排序损失、边界 shortcut 稀疏化、动态需求更新和
更多城市复现。对于 Chicago 一类密集边界图，优先优化扫边预算可能比继续提高全局 Spearman 更
直接。

\section{有效性威胁}

\paragraph{内部有效性。}
区域标签来自精确查询工作量，但展开节点仍受双向 Dijkstra 的终止条件与等长路径 tie 顺序影响。
为降低实现差异，所有同组比较共享算法和数据结构；C++ 与 Python 的平均展开差异被限制在每查询
0.1 个节点以内。训练过程曾经历结构和损失修正，但最终 \gmethod 协议在解锁 holdout 前冻结，
旧失败运行完整保留，不进入正式汇总。

\paragraph{构念有效性。}
Spearman 衡量全局排序，NDCG 与 TopGain 衡量头部，展开节点和扫描边衡量算法工作量，墙钟时间
衡量具体实现效果。这些指标不能相互替代。尤其是区域标签采用单区域展开节点收益，不直接表示
多区域加速；本文以 hard-disjoint 在线配对补足这一缺口。公平性、用户体验和交通可达性未被直接
测量，因此不作相关结论。

\paragraph{外部有效性。}
Porto 出租车与 Chicago 共享单车提供了跨城市、跨出行方式证据，但仍只有两个城市，且都使用
固定区域尺度和静态长度边权。结论更适用于具有重复 OD 查询、可离线预处理且要求距离精确的
道路分析任务。对于实时旅行时间、频繁闭路、负权或强动态拓扑，需重新评估传播输入与索引更新
成本。

\paragraph{结论有效性。}
\gmethod 的三个种子全部在两类外部窗口上提高 Spearman，支持方向稳定性，但样本数不足以声称
普遍统计优势；头部指标中的明显反例已经保留。Chicago 平均耗时未改善也未被描述为失败，而是
用于定位 shortcut 密度和扫边成本这一适用边界。

\section{结论}

本文在已有 CRP 式精确边界覆盖图上研究历史 OD 驱动的选择性区域物化。通过无路径输入的
\zmethod、以其为冻结先验的 \gmethod，以及时间、空间和部署三层隔离协议，本文把离线区域排序
与精确在线最短路查询连接为可审计闭环。两城结果表明，正确道路拓扑上的历史边际需求传播能够
提供稳定且经济的区域价值排序；神经残差可进一步改善全局排序，但其头部与系统增量受预算、
shortcut 密度和图结构限制。该结论既保留了 AI 的可验证增量，也避免以近似模型牺牲最短路答案
正确性。

\appendix

\section{冻结配置摘要}

\begin{table}[H]
  \centering
  \caption{正式实验的关键冻结配置。}
  \begin{tabularx}{\textwidth}{lX}
    \toprule
    项目 & 固定值 \\
    \midrule
    时间窗口 & $H=[0,0.35)$，$Y=[0.35,0.70)$，$F=[0.70,1.00)$ \\
    候选 & 每城 1,200 个，每候选 512 节点，候选种子 42 \\
    空间隔离 & Jaccard 阈值 0.50 的重叠连通组，840/180/180 \\
    标签查询 & 每城、每窗口、每候选 2,000 条，抽样种子 42，端点缓存关闭 \\
    \zmethod & 正反向 lazy mean，深度 1/2/4/8/16/32，区域 mean+max \\
    \gmethod & 32 层双塔，hidden dim 32，残差尺度 0.01，418,183 参数 \\
    优化 & CUDA FP32，学习率 $5\times10^{-3}$，weight decay 0，40 epoch，无 scheduler \\
    模型选择 & seeds 42/43/44，validation Spearman，$\alpha\in\{0,0.125,0.25,0.5,1\}$ \\
    部署 & hard-disjoint，$K=5/10/18$；主系统表报告 $K=18$ \\
    C++ 计时 & C++20/O3，固定单核，2 轮预热，10 轮正式重复，每组 2,000 查询 \\
    \bottomrule
  \end{tabularx}
\end{table}

\section{论文主张边界}

为保持相关工作归因准确，本文不主张发明区域受限 Dijkstra、边界 clique shortcut 或双向图扩散
基础算子，也不主张首次利用历史查询负载优化精确最短路索引。本文可复现实证支持的范围是：在
既有精确覆盖图底座上，历史 OD 驱动的选择性区域物化可以通过无路径的确定性需求传播实现；以
该方法为先验的神经残差可稳定改善全局排序；时空隔离、非重叠选择和在线配对能够揭示离线相关
性无法表达的成本与适用边界。

\section*{声明}

\noindent 数据可用性：\rule{0.76\linewidth}{0.4pt}\\[0.8em]
代码可用性：\rule{0.76\linewidth}{0.4pt}\\[0.8em]
基金项目：\rule{0.79\linewidth}{0.4pt}\\[0.8em]
利益冲突：\rule{0.79\linewidth}{0.4pt}\\[0.8em]
作者贡献：\rule{0.79\linewidth}{0.4pt}\\[0.8em]
致谢：\rule{0.84\linewidth}{0.4pt}

\bibliographystyle{plainnat}
\bibliography{references}

\end{document}
